\subsection{Contact transformations.}

In the following two years, Lie seems to abandon the theory of transformation groups (although he remains in very close contact with Klein, who publishes in 1872 his ``Programme'') in order to study contact transformations, the integration of partial differential equations of the first order and the relations between these two theories. We do not have to spell out the history of these questions here, and we will limit ourselves to mentioning a few points that seem to have played an important role in the genesis of the theory of transformation groups.

The notion of contact transformation generalises at the same time the point transformations and the transformations by reciprocal polars. Grosso modo, a contact transformation \(^{1}\) in \(C^{n}\) is an isomorphism of an open set \(\Omega\) of the variety \(T^{\prime}(C^{n})\) of cotangent vectors to \(C^{n}\) onto another open set \(\Omega^{\prime}\) of \(T^{\prime}(C^{n})\) transforming the canonical 1-form of \(\Omega\) into that of \(\Omega^{\prime}\) . In other words, if \((x_{1},\ldots,x_{n},p_{1},\ldots,p_{n})\) describes the canonical co-ordinates of \(T^{\prime}(C^{n})\) , a contact transformation is an isomorphism \((x_{i},p_{i})\mapsto(X_{i},P_{i})\) satisfying the relation \(\sum_{i=1}^{n}P_{i}dX_{i}=\sum_{i=1}^{n}p_{i}dx_{i}\) . Such transformations occur in the study of partial differential equations of the form

\[
F \left(x _ {1}, x _ {2}, \dots , x _ {n}, \frac {\partial z}{\partial x _ {1}}, \dots , \frac {\partial z}{\partial x _ {n}}\right) = 0\tag{25.2}
\]

Lie becomes familiar during the course of his research on these questions with the handling of Poisson parentheses

\[
(f, g) = \sum_ {i = 1} ^ {n} \left(\frac {\partial f}{\partial x _ {i}} \frac {\partial g}{\partial p _ {i}} - \frac {\partial g}{\partial x _ {i}} \frac {\partial f}{\partial p _ {i}}\right)\tag{25.3}
\]

and brackets \(^{2}\) \([X,Y]=XY-YX\) of differential operators of type (25.1); he interprets the Poisson parenthesis (25.3) as the effect on f of a transformation of type (25.1) associated with g, and observes on this occasion that the Jacobi identity for Poisson parentheses means that the bracket of differential operators corresponding to g and h is associated with the parenthesis \((g,h)\) . The search for functions g such that \((F,g)=0\) , which occurs in the method of Jacobi for integrating the partial differential equation (25.2), becomes for Lie that for an infinitesimal contact transformation leaving the given equation invariant. Finally, Lie is led to study the set of functions \((u_{j})_{1\leq j\leq m}\) of the \(x_{i}\) and the \(p_{i}\) such that the parentheses \((u_{j},u_{k})\) are functions of the \(u_{h}\) , and calls these sets ``groups'' (already considered in substance by Jacobi).

